\documentclass[10pt,twocolumn]{article}

\usepackage[letterpaper,margin=0.75in]{geometry}
\usepackage{amsmath,amssymb}
\usepackage{booktabs}
\usepackage{graphicx}
\usepackage{microtype}
\usepackage{xcolor}
\usepackage{enumitem}
\usepackage[hidelinks]{hyperref}
\usepackage[T1]{fontenc}
\usepackage{lmodern}

\setlength{\columnsep}{0.24in}
\setlength{\parindent}{1em}
\setlength{\parskip}{0pt}
\setlist{nosep,leftmargin=*}

\newcommand{\method}{\textsc{BilingualAnchor}}
\newcommand{\todo}[1]{\textcolor{red}{[TODO: #1]}}
\newcommand{\dataset}{\textsc{Ko-JB}}

\title{Guarding Korean without Korean Safety Labels:\\
Middle-Layer Alignment for Cross-Lingual Jailbreak Detection}

\author{Anonymous Authors \\
Anonymous Institution \\
\texttt{anonymous@example.com}}

\date{}

\begin{document}
\maketitle

\begin{abstract}
Input guard models are commonly trained with English safety data, but their
decisions may transfer poorly to languages that are absent from task-specific
supervision. We study Korean jailbreak detection under a constrained setting:
English prompts and teacher predictions are available, together with a small
set of English--Korean translation pairs, but no Korean safety labels are used
for training. We introduce \method, which jointly distills the task decision
of Llama Prompt Guard 2 into a multilingual DeBERTa encoder and contrastively
aligns mean-pooled English and Korean representations at an intermediate
layer. We evaluate on a parallel English--Korean benchmark of 2,485 test
prompts, including 1,848 attacks from temporally held-out,
in-the-wild, and synthetic adversarial sources. At a validation-calibrated 1\%
false-positive-rate operating point, alignment at layer 8 raises Korean recall
from 24.4\% to 33.2\% over a matched distillation baseline while leaving
English recall essentially unchanged (33.5\% versus 33.4\%). The gain is
largest on WildJailbreak (3.7\% to 14.2\%), and a paired McNemar test finds a
significant difference ($p<10^{-34}$). Layer-wise ablations show a non-monotonic
effect: shallow alignment hurts transfer, whereas middle-to-upper-layer
alignment is consistently beneficial. These results suggest that unlabeled
bilingual anchoring can improve target-language guard transfer, while also
exposing substantial remaining challenges in strict low-FPR calibration and
out-of-distribution recall.
\end{abstract}

\section{Introduction}

Large language model (LLM) applications increasingly rely on input-side guard
models to reject jailbreaks and prompt injections before a user request reaches
the generator. Such guards must operate at low false-positive rates (FPRs): a
detector that rejects a large fraction of benign traffic is difficult to deploy,
even if its aggregate classification accuracy is high. Multilingual deployment
adds a second requirement. A model trained primarily on English attack patterns
must preserve its decision boundary when equivalent instructions are expressed
in another language.

Korean is a useful stress test for this transfer problem. It is typologically
distant from English, uses a different writing system, and is not among the
eight evaluation languages reported for Llama Prompt Guard 2
\cite{meta2025promptguard2}. Collecting Korean safety annotations for every new
attack distribution is expensive and slow, particularly because jailbreak
prompts evolve over time \cite{piet2025jot}. We therefore ask:

\begin{quote}
Can an English-trained input guard acquire stronger Korean jailbreak detection
without using Korean safety labels?
\end{quote}

We investigate this question by separating task transfer from language
alignment. The task objective distills soft predictions from Llama Prompt Guard
2 on 18,585 English prompts. A second objective takes 1,958 English--Korean
translation pairs and makes corresponding representations similar at a chosen
student layer. The alignment examples carry no task labels in the loss. This
design aims to preserve an English teacher's attack knowledge while repairing
the bilingual geometry of a multilingual encoder.

Our experiments yield three main findings. First, the distillation-only student
transfers unevenly: it detects most temporally held-out Korean JOT attacks but
only 3.7\% of Korean WildJailbreak examples at the chosen operating point.
Second, alignment at layer 8 improves overall Korean recall by 8.8 percentage
points without a measurable loss in English recall relative to the matched
baseline. Third, the effect depends strongly on depth: layer 4 alignment is
harmful, layers 6--10 help, and layer 8 provides the best recall on our diverse
main benchmark. The result supports intermediate representation alignment as a
practical target-language adaptation mechanism, but the low absolute recall on
the hardest source also cautions against interpreting cross-lingual improvement
as solved robustness.

Our contributions are:
\begin{itemize}
    \item an adaptation objective that combines English-only task distillation
    with label-free English--Korean representation alignment for an encoder
    guard model;
    \item a source-diverse, exactly parallel English--Korean evaluation suite
    with explicit low-FPR calibration; and
    \item a layer- and source-level analysis showing where bilingual alignment
    helps, where it hurts, and how representation retrieval changes after
    training.
\end{itemize}

\paragraph{Terminology.}
Our setting is \emph{zero Korean safety labels}, not zero-shot Korean: Korean
translations are used by the alignment objective. We also use ``distillation''
to mean teacher-to-student knowledge transfer, not model compression; the
student checkpoint is larger than the 86M-parameter teacher.

\section{Related Work}

\subsection{Jailbreak detection and evaluation}

In-the-wild jailbreaks use role play, instruction override, privilege
escalation, and other social or structural devices to bypass model policies
\cite{shen2023dan}. WildTeaming mines such tactics and composes them to create
the large-scale WildJailbreak resource \cite{jiang2024wildteaming}.
JailbreaksOverTime (JOT) shows that attack distributions drift and provides
timestamped benign and attack prompts for temporal evaluation
\cite{piet2025jot}. These findings motivate our use of both a temporal holdout
and sources not used for task distillation.

Prompt Guard 2 is a binary DeBERTa-family classifier for explicit jailbreak and
prompt-injection patterns \cite{meta2025promptguard2}. Guard evaluation must
also account for over-defense: benign prompts containing attack-adjacent words
can be incorrectly rejected. NotInject was introduced specifically to measure
this behavior \cite{li2025piguard}, and we include its translated benign prompts
in calibration and testing.

\subsection{Cross-lingual knowledge transfer}

Knowledge distillation can transfer task predictions from a monolingual teacher
to a multilingual student without target-language task labels.
MONOX-KD applies this idea to cross-lingual classification
\cite{chi2020monox}. Distillation alone, however, does not explicitly constrain
the relation between parallel source- and target-language representations.

Liu and Niehues \cite{liu2025middle} show that middle-layer representation
alignment improves cross-lingual transfer in fine-tuned decoder-only LLMs. We
adapt this principle to binary encoder-based guard classification. Our
implementation differs in two respects: it jointly optimizes task and alignment
losses rather than alternating between them, and it uses in-batch contrastive
alignment of mean-pooled sentence representations. Accordingly, our novelty is
the guard adaptation setting and its Korean low-FPR analysis, not the general
idea of middle-layer alignment.

\section{Problem Setting}

Let $x^{en}$ be an English prompt and $y\in\{0,1\}$ indicate benign or
malicious input. A frozen teacher $T$ provides a distribution
$q_T(y\mid x^{en})$. The student $S_\theta$ is a multilingual encoder classifier.
During training, we have an English task corpus
$\mathcal{D}_{task}=\{x_i^{en}\}$ and a separate parallel corpus
$\mathcal{D}_{par}=\{(a_j^{en},a_j^{ko})\}$. Safety labels associated with the
parallel corpus are not consumed by the alignment loss.

The goal is to improve Korean detection while retaining English behavior. At
evaluation time, a threshold is selected only from a validation set of benign
prompts to target a specified FPR and is then frozen for the test set.

\section{Method}

\subsection{English task distillation}

For each English task prompt, we minimize cross-entropy against the teacher's
soft distribution:
\begin{equation}
\mathcal{L}_{task}
=-\frac{1}{B}\sum_{i=1}^{B}\sum_{c\in\{0,1\}}
q_T(c\mid x_i^{en})\log p_\theta(c\mid x_i^{en}).
\end{equation}
No ground-truth task label is used in this objective. Both distributions use
temperature 1.

\subsection{Layer-specific bilingual alignment}

Let $H_\theta^\ell(x)\in\mathbb{R}^{m\times d}$ be token representations at
layer $\ell$. We mean-pool non-padding tokens and L2-normalize the result:
\begin{equation}
z_\ell(x)=\operatorname{norm}_2\left(
\frac{\sum_t M_t H_{\theta,t}^\ell(x)}{\sum_t M_t}\right).
\end{equation}
For a minibatch of $B$ translation pairs, the one-directional in-batch
contrastive loss is
\begin{equation}
\mathcal{L}_{align}^{\ell}
=-\frac{1}{B}\sum_{i=1}^{B}\log
\frac{\exp(z_\ell(a_i^{en})^\top z_\ell(a_i^{ko})/\tau)}
{\sum_{j=1}^{B}\exp(z_\ell(a_i^{en})^\top z_\ell(a_j^{ko})/\tau)},
\end{equation}
where $\tau=0.05$. Other Korean sentences in the batch act as negatives.

The final objective is
\begin{equation}
\mathcal{L}=\mathcal{L}_{task}+\lambda\mathcal{L}_{align}^{\ell},
\qquad \lambda=1.
\end{equation}
We train variants with $\ell\in\{4,6,8,10\}$. The matched L0 baseline uses the
same task pipeline but omits the parallel inputs and alignment loss.

\section{Data}

\subsection{Training data}

The task corpus contains 18,585 unique English prompts from JOT between
February and November 2023. It contains 3,265 attacks and 15,320 benign prompts,
although these labels are retained only for corpus analysis and are not used by
the distillation objective. Exact prompt overlap with the December 2023 split is
removed before training.

The alignment corpus contains 1,958 English--Korean machine-translated pairs
drawn from the English training pool (346 attacks and 1,612 benign prompts by
the unused reference labels). All English strings, Korean strings, and pair
identifiers are unique, and we find no exact prompt overlap with the evaluation
sets. Because the task corpus is larger, alignment pairs are traversed
cyclically while iterating over task examples.

\subsection{Evaluation data}

Table~\ref{tab:data} summarizes the datasets in scope. The main \dataset{}
test set is exactly parallel: each English prompt has a Korean translation with
the same identifier and label. Both attack and benign classes are translated,
reducing a trivial translation-style/class correlation. The 1,848 attacks
combine temporal JOT attacks, prompts collected in the wild, and WildJailbreak.
The 637 benign examples combine JOT and NotInject. The validation split contains
300 benign prompts and is used only for threshold calibration.

\begin{table}[t]
\centering
\small
\begin{tabular}{lrrr}
\toprule
Split & Attack & Benign & Total \\
\midrule
JOT train (English) & 3,265 & 15,320 & 18,585 \\
Alignment pairs & 346 & 1,612 & 1,958 \\
\midrule
\dataset{} validation & 0 & 300 & 300 \\
\dataset{} test & 1,848 & 637 & 2,485 \\
JOT-Dec parallel test & 286 & 682 & 968 \\
\bottomrule
\end{tabular}
\caption{Data statistics. Reference labels in the alignment rows are not used
by the alignment objective. All evaluation counts apply to each language.}
\label{tab:data}
\end{table}

The attack composition of the main test set is 1,168 WildJailbreak examples,
394 in-the-wild examples, and 286 JOT December examples. Benign examples contain
467 JOT and 170 NotInject prompts.

\paragraph{Translation status.}
The December JOT translations have mean multilingual embedding similarity
0.873 (median 0.886). Automated screening flags 41 broader benchmark examples
for manual review. \todo{Complete the 41-item intent-preservation review,
rebuild the benchmark, and update all numbers before submission. Specify the
translation model/API, version, decoding settings, and human-review protocol.}

\section{Experimental Setup}

We initialize the student from \texttt{microsoft/mdeberta-v3-base}, a 12-layer
multilingual DeBERTa V3 encoder \cite{he2021debertav3}, and the classifier head
for two labels. The teacher is
\texttt{meta-llama/Llama-Prompt-Guard-2-86M}. The student has approximately 279M
parameters due largely to its 250K multilingual vocabulary; its saved weights
occupy 1.1GB.

All current models are trained for three epochs with learning rate
$2\times10^{-5}$, batch size 16, 10\% warmup, weight decay 0.01, FP16, maximum
sequence length 512, and seed 0. Training uses no validation-based checkpoint
selection. \todo{Repeat L0 and L8 with at least three, preferably five, seeds and
report mean, standard deviation, and per-seed operating points.}

For each model and language, we score maliciousness as $1-p_\theta(y=0\mid x)$.
We sort benign validation scores and select the lowest threshold whose empirical
validation FPR does not exceed the target $\alpha=0.01$. The threshold is then
applied unchanged to the test set. We report both recall and the achieved test
FPR. We use an exact two-sided McNemar test on paired attack decisions for
significance. Because all current trained results use one seed, this test
captures evaluation-item uncertainty but not training variance.

\section{Results}

\subsection{Main bilingual evaluation}

Table~\ref{tab:main} reports the main result. L8 reaches 33.2\% Korean recall,
compared with 24.4\% for the matched L0 baseline, an 8.8-point absolute and
36.1\% relative improvement. English recall is effectively unchanged
(33.4\% versus 33.5\%). Thus the improvement is not obtained by exchanging
English detection for Korean detection.

\begin{table*}[t]
\centering
\small
\begin{tabular}{lrrrrr}
\toprule
Model & Align layer & English recall & English test FPR & Korean recall & Korean test FPR \\
\midrule
Legacy MONOX-KD & -- & 35.7 & 1.10 & 24.3 & 2.35 \\
Matched KD (L0) & -- & 33.5 & 1.73 & 24.4 & 1.41 \\
\method{} & 4 & 28.1 & 2.04 & 20.1 & 0.78 \\
\method{} & 6 & 35.5 & 2.04 & 30.2 & 1.41 \\
\method{} & 8 & 33.4 & 1.41 & \textbf{33.2} & 1.57 \\
\method{} & 10 & 34.8 & 2.67 & 32.5 & 1.88 \\
Prompt Guard 2 teacher & -- & \textbf{47.6} & 1.73 & 27.1 & 0.78 \\
\bottomrule
\end{tabular}
\caption{Recall and achieved test FPR (\%) on the main parallel \dataset{} test
set. Thresholds are chosen independently for each language and model using a
1\% target FPR on its validation split. These are validation-calibrated
operating points; achieved test FPR need not equal 1\%.}
\label{tab:main}
\end{table*}

The teacher remains substantially stronger in English, exceeding L8 by 14.2
points. L8 exceeds the teacher's Korean recall by 6.1 points, but this is not a
strict equal-FPR comparison because the achieved Korean test FPR is 1.57\% for
L8 and 0.78\% for the teacher. We therefore use the matched L0 comparison as our
primary causal evidence for alignment.

The paired comparison further supports the main effect. Among 1,848 Korean
attacks, 181 are missed by L0 but detected by L8, while 18 show the reverse
pattern. The resulting exact McNemar $p$-value is $4.68\times 10^{-35}$.
Layer 6 and layer 10 also significantly improve over L0, whereas layer 4
significantly degrades performance. \todo{After multi-seed training, supplement
paired item tests with uncertainty over seeds and correct for multiple layer
comparisons.}

\subsection{Where does alignment help?}

Table~\ref{tab:source} decomposes Korean attack recall by source. The largest
absolute gain occurs on WildJailbreak: L8 improves from 3.7\% to 14.2\%. JOT is
the easiest source for all models, likely because it shares the data-generating
family of the distillation corpus despite the temporal split. In-the-wild
prompts occupy the middle. The source decomposition shows that the aggregate
gain is not solely driven by the temporally held-out JOT subset.

\begin{table}[t]
\centering
\small
\begin{tabular}{lrrr}
\toprule
Source & L0 & L8 & $\Delta$ \\
\midrule
JOT Dec. & 82.2 & 89.9 & +7.7 \\
In-the-wild & 44.2 & 48.5 & +4.3 \\
WildJailbreak & 3.7 & 14.2 & +10.5 \\
\bottomrule
\end{tabular}
\caption{Korean attack recall (\%) by source at each model's
validation-calibrated 1\% target-FPR operating point.}
\label{tab:source}
\end{table}

The absolute WildJailbreak recall remains low even after alignment. Our method
should therefore be understood as improving cross-lingual transfer under source
shift, not as solving out-of-distribution jailbreak detection.

\section{Analysis}

\subsection{Layer choice is non-monotonic}

Alignment is not uniformly helpful. L4 reduces Korean recall by 4.3 points,
whereas L6, L8, and L10 improve it by 5.8, 8.8, and 8.1 points, respectively.
This pattern is consistent with the hypothesis that shallow representations
must retain language- and form-specific information, while intermediate
representations offer a better interface for aligning task-relevant semantics.
The slight decline from L8 to L10 suggests that constraining representations
closer to the classifier is not always optimal.

\subsection{Parallel retrieval}

We evaluate representation alignment without safety labels by retrieving the
English counterpart of each Korean sentence among 600 candidates. Table
\ref{tab:retrieval} reports top-1 accuracy at student layers 8 and 12. Training
at L8 gives the highest layer-8 retrieval accuracy, improving it from 57.8\% to
71.7\%. Training at L10 produces the highest final-layer retrieval but lower
main-benchmark recall than L8. Thus stronger alignment everywhere is not
sufficient; where bilingual geometry is imposed matters for downstream
transfer.

\begin{table}[t]
\centering
\small
\begin{tabular}{lrr}
\toprule
Training variant & Retrieval at L8 & Retrieval at L12 \\
\midrule
L0 & 57.8 & 9.3 \\
L4 & 55.8 & 8.3 \\
L6 & 61.7 & 14.5 \\
L8 & \textbf{71.7} & 22.5 \\
L10 & 67.0 & \textbf{43.5} \\
\bottomrule
\end{tabular}
\caption{English--Korean top-1 parallel retrieval accuracy (\%) among 600
candidates, measured at two representation layers.}
\label{tab:retrieval}
\end{table}

\subsection{Calibration at low FPR}

The validation-calibrated target does not guarantee the same test FPR under
distribution shift. For example, L8 targets 1\% but reaches 1.57\% on Korean
test data. More fundamentally, 300 benign validation examples provide only
three-error resolution at 1\% FPR. A 0.1\% target would allow zero errors and is
not statistically resolvable with this validation size. We consequently avoid
using 0.1\% as a headline metric and report achieved FPR beside recall.

\todo{Before submission, enlarge the benign calibration set; store per-example
scores; report bootstrap confidence intervals, ROC curves, and partial AUC in
the low-FPR region; and verify numerical stability around tied saturated
scores.}

\section{Limitations}

\paragraph{Single language pair.}
We study English-to-Korean transfer only. The result does not establish that
the same layer is optimal for other languages or scripts.

\paragraph{Machine translation.}
Alignment and evaluation data depend on machine translation. Automated semantic
similarity cannot guarantee that attack intent, ambiguity, or culturally
specific wording is preserved. The current draft precedes completion of a
41-item flagged translation review. All final results must be recomputed after
that audit.

\paragraph{Layer selection.}
Layer 8 is best on the main benchmark, whereas layer 10 is slightly stronger on
the JOT-only parallel evaluation under the current calibration. Unless an
independent, pre-specified layer selection rule is documented, L8 should be
reported as the best observed layer in a sweep rather than an a priori selected
architecture.

\paragraph{Calibration sample size.}
The benign validation split is too small for precise estimation below 1\% FPR.
Differences between validation-targeted and achieved test FPR complicate direct
comparison between independently calibrated models.

\paragraph{Training variance and controls.}
Current results use one seed. We have not yet evaluated shuffled translation
pairs, Korean pseudo-label distillation, or translate-test baselines. Such
controls are needed to distinguish semantic alignment from additional
regularization and compute.

\paragraph{Scope.}
The method is evaluated as a jailbreak detector. Preliminary evaluation on the
PIArena indirect prompt-injection task yields near-zero recall for the aligned
models, indicating a task mismatch. We do not claim general prompt-injection or
agent-security robustness.

\paragraph{Efficiency.}
The student is larger than the teacher and alignment adds two encoder forward
passes during training. This work addresses target-language transfer rather
than compression or inference efficiency.

\section{Ethical Considerations}

The benchmark contains prompts requesting harmful, illegal, hateful, or abusive
content. Annotators should be informed of this material and provided with the
ability to stop review. Released data should preserve source licenses and avoid
including personally identifying information. Because the classifier can both
miss attacks and reject legitimate requests, it should be deployed as one layer
of a defense-in-depth system rather than as a sole safety mechanism. Reporting
source-level failures and achieved false-positive rates is important to avoid a
misleading impression of protection.

\section{Reproducibility Status}

The training code fixes the seed, model identifiers, maximum length, loss
weights, temperature, and optimizer hyperparameters. The temporal JOT split and
exact-overlap removal are scripted. The current repository also stores final
model checkpoints and paired hit vectors.

The following artifacts remain necessary for a reproducible release:
\begin{itemize}
    \item a deterministic script that builds \dataset{} and documents the
    exclusion of 46 translated source prompts;
    \item reviewed and versioned translation files with provenance and licenses;
    \item per-example validation and test scores, thresholds, and achieved FPRs;
    \item multi-seed checkpoints and aggregate statistics; and
    \item exact-benchmark outputs for external multilingual guard baselines.
\end{itemize}

\section{Conclusion}

We presented \method, a simple combination of English teacher distillation and
label-free bilingual representation alignment for Korean jailbreak detection.
The best observed intermediate-layer variant improves Korean recall by 8.8
points over matched distillation while preserving English recall, with the
largest gain on attacks outside the task-distillation source. The non-monotonic
layer effect and retrieval analysis indicate that cross-lingual guard transfer
depends not only on whether representations are aligned but also on where the
constraint is applied. At the same time, low absolute recall on WildJailbreak
and calibration uncertainty at strict FPRs show that bilingual alignment is a
useful adaptation tool rather than a complete defense.

\section*{Acknowledgments}

\todo{Add funding, compute, translation-review, and institutional acknowledgments.}

\begin{thebibliography}{99}

\bibitem[Chi et~al.(2020)Chi, Dong, Wei, Mao, and Huang]{chi2020monox}
Zewen Chi, Li Dong, Furu Wei, Xianling Mao, and Heyan Huang. 2020.
\newblock Can monolingual pretrained models help cross-lingual classification?
\newblock In \emph{Proceedings of AACL-IJCNLP}, pages 12--17.

\bibitem[He et~al.(2023)He, Gao, and Chen]{he2021debertav3}
Pengcheng He, Jianfeng Gao, and Weizhu Chen. 2023.
\newblock DeBERTaV3: Improving DeBERTa using ELECTRA-style pre-training with
gradient-disentangled embedding sharing.
\newblock In \emph{Proceedings of ICLR}.

\bibitem[Jiang et~al.(2024)Jiang, Rao, Han, Ettinger, Brahman, Kumar,
Mireshghallah, Lu, Sap, Choi, and Dziri]{jiang2024wildteaming}
Liwei Jiang, Kavel Rao, Seungju Han, Allyson Ettinger, Faeze Brahman,
Sachin Kumar, Niloofar Mireshghallah, Ximing Lu, Maarten Sap, Yejin Choi,
and Nouha Dziri. 2024.
\newblock WildTeaming at scale: From in-the-wild jailbreaks to
(adversarially) safer language models.
\newblock In \emph{Advances in Neural Information Processing Systems}.

\bibitem[Li et~al.(2025)Li, Liu, Zhang, and Xiao]{li2025piguard}
Hao Li, Xiaogeng Liu, Ning Zhang, and Chaowei Xiao. 2025.
\newblock PIGuard: Prompt injection guardrail via mitigating overdefense for free.
\newblock In \emph{Proceedings of ACL}, pages 30420--30437.

\bibitem[Liu and Niehues(2025)]{liu2025middle}
Danni Liu and Jan Niehues. 2025.
\newblock Middle-layer representation alignment for cross-lingual transfer in
fine-tuned LLMs.
\newblock In \emph{Proceedings of ACL}, pages 15979--15996.

\bibitem[Meta(2025)]{meta2025promptguard2}
Meta. 2025.
\newblock Llama Prompt Guard 2 model card.
\newblock \url{https://huggingface.co/meta-llama/Llama-Prompt-Guard-2-86M}.

\bibitem[van~den Oord et~al.(2018)van~den Oord, Li, and Vinyals]{oord2018cpc}
Aaron van~den Oord, Yazhe Li, and Oriol Vinyals. 2018.
\newblock Representation learning with contrastive predictive coding.
\newblock \emph{arXiv preprint arXiv:1807.03748}.

\bibitem[Piet et~al.(2025)Piet, Huang, Jacob, Chow, Alrashed, Zhao, Hu,
Sitawarin, Alomair, and Wagner]{piet2025jot}
Julien Piet, Xiao Huang, Dennis Jacob, Annabella Chow, Maha Alrashed,
Geng Zhao, Zhanhao Hu, Chawin Sitawarin, Basel Alomair, and David Wagner. 2025.
\newblock JailbreaksOverTime: Detecting jailbreak attacks under distribution
shift.
\newblock In \emph{Proceedings of the ACM Workshop on Artificial Intelligence
and Security}.

\bibitem[Shen et~al.(2023)Shen, Chen, Backes, Shen, and Zhang]{shen2023dan}
Xinyue Shen, Zeyuan Chen, Michael Backes, Yun Shen, and Yang Zhang. 2023.
\newblock ``Do Anything Now'': Characterizing and evaluating in-the-wild
jailbreak prompts on large language models.
\newblock \emph{arXiv preprint arXiv:2308.03825}.

\end{thebibliography}

\appendix

\section{Implementation Details}

Teacher scores are computed once over English task prompts and stored in memory
for training. The student receives one task batch and, for aligned variants,
two additional forward passes for the English and Korean members of a parallel
batch. Padding tokens are excluded from mean pooling. The alignment target is
the diagonal of the $B\times B$ similarity matrix. Models are saved only after
the third epoch; there is no best-checkpoint selection despite the checkpoint
directory name \texttt{best}.

\section{Additional Statistical Results}

Using the stored paired Korean attack decisions, the exact McNemar comparisons
against L0 are shown in Table~\ref{tab:mcnemar}. ``Gain'' counts attacks missed
by L0 and detected by the aligned model; ``loss'' is the reverse.

\begin{table}[h]
\centering
\small
\begin{tabular}{lrrr}
\toprule
Variant & Gain & Loss & $p$ \\
\midrule
L4 & 27 & 107 & $1.82\times10^{-12}$ \\
L6 & 134 & 26 & $9.92\times10^{-19}$ \\
L8 & 181 & 18 & $4.68\times10^{-35}$ \\
L10 & 178 & 28 & $6.78\times10^{-28}$ \\
\bottomrule
\end{tabular}
\caption{Exact paired McNemar tests on 1,848 Korean attacks. These tests do not
include training-seed uncertainty.}
\label{tab:mcnemar}
\end{table}

\section{Planned Controls}

The camera-ready experimental matrix should add the following controls while
keeping the main test set frozen:
\begin{enumerate}
    \item Korean-to-English translate-test followed by the frozen teacher;
    \item pseudo-label KD on the same 1,958 Korean translations;
    \item alignment with randomly permuted English--Korean pairs;
    \item a matched-compute regularization control; and
    \item external multilingual guards evaluated on the exact same examples.
\end{enumerate}

\end{document}
